\section{Supporting protocol status before the baseline}
\label{sec:status}

This manuscript is the supporting protocol record; the results report will be
the primary research artifact once scored runs exist. The design above is
frozen, but the experiment has not been run. Stating the exact boundary matters,
because a protocol paper that blurs it invites its own hypotheses to be read as
findings.

Implemented and committed: the reproducible preparation of the eligible manifest
from the pinned public release; both deterministic splits with byte-identical
regeneration tests; the custody release tool and its verification against the
externally recorded pre-gold commit; both scorers as pure comparators with
synthetic conformance tests; the append-only ledger, experiment lifecycle,
checkpoint metering, and receipt verification of the development control plane;
the generation and score artifact schemas with recursive rejection of protected
fields; and a credential-safe contract probe for the production system that
verifies configuration and manifests against the recorded commit before
authenticating, forwards only minimal authentication environment, projects
responses into type and shape metadata, hashes response bodies, and writes only
private artifacts under ignored roots. The sealed PostgreSQL execution layer is
implemented and conformance-tested. All 18 public database mirrors have passed
schema and content parity plus read-only-role verification. C1 through C3 now
have committed condition configurations and a traceable direct-SQL runtime.

One public-only governed-product canary has also run on an isolated model branch.
Its semantic bundle passed product validation and exact semantic readback; one
read-only governed query and one unscored AI Hub diagnostic completed. This was
an integration check, not a benchmark question or correctness judgment.

Not done: no sealed generation exists, no output has been scored, and no accuracy
of any kind has been computed for any condition. The capture verification gate of
Section~\ref{sec:conditions} has not been passed. C1 through C3 still lack the
executable driver needed for a live attempt, and no condition has yet emitted the
complete scaled-baseline envelope. The public-only semantic bundle covers one of
18 databases and only the mechanically supported portion of that database's HKB.
The guardian public-key hash is pinned in configuration while the private key
remains outside the development workspace. Exact final model identifiers, retry
and time ceilings, version fingerprints, immutable C3/C4 model identity, and the
C4 result-set adaptation path remain pre-baseline or Freeze B work.

The failure taxonomy currently contains risk hypotheses, not observations. The
three we consider most likely, stated here so they can be scored against what is
eventually measured, are HKB dependency composition (a mechanical compiler may
represent nodes while losing dependency meaning, grain, or composition), semantic
discoverability (a definition may be correct in the model yet absent from the
retrieval surface the agent uses, producing an apparent reasoning failure that is
really context selection), and compilation and validation fidelity (correct
semantic intent may be altered during compilation or rejected by validation and
retry behavior). Every taxonomy cell is currently marked unmeasured, and the
highest-information next experiment is the public-only mechanical baseline with
trace capture sufficient to classify each knowledge-linked development failure on
the mechanism ladder of Section~\ref{sec:analysis}.

A results paper will follow the sealed event. It will report the public-only
mechanical baseline, the supervised experiment history including regressions and
dead ends, the leave-one-database-out portability results on development data,
the frozen held-out primary result and condition matrix, and the reliability,
failure, cost, latency, and product analyses, each derivable from preserved
machine-readable artifacts. Whatever it reports, the design it reports against is
the one in this paper.
